\section{Categorías sobre otra categoría}
\label{VI.2}

Sea $\cal{E}$ una categoría en~$\Univ$; es, por tanto, un objeto de~$\Cat$, y
se puede considerar la categoría
$\Cat_{/\cal{E}}$
\label{indnot:fh}\oldindexnot{$\Cat_{/\cal{E}}$|hyperpage}%
de los «objetos de~$\Cat$ sobre~$\cal{E}$». Un objeto de esta categoría es,
pues, un funtor
$$
p\colon \cal{F}\to\cal{E}
$$
También se dice que la categoría~$\cal{F}$, provista de tal funtor, es una
\emph{categoría sobre}~$\cal{E}$, o una $\cal{E}$\emph{-categoría}. Se llamará,
por tanto, $\cal{E}$-funtor de una categoría~$\cal{F}$ sobre~$\cal{E}$ en una
categoría~$\cal{G}$ sobre~$\cal{E}$ a un funtor
$$
f\colon\cal{F}\to\cal{G}
$$
tal que
$$
qf=p
$$
donde $p$ y $q$ son los funtores de proyección de~$\cal{F}$ y~$\cal{G}$,
respectivamente. El conjunto de los $\cal{E}$-funtores~$f$ de~$\cal{F}$ en
$\cal{G}$ está, por tanto, en correspondencia biunívoca con el conjunto de las
flechas de origen~$\cal{F}$ y extremo~$\cal{G}$ en~$\Cat_{/\cal{E}}$, sin
\marginpar{149}
que se trate por ello de una identidad (pues el dato de un~$f$ como el anterior
no determina $\cal{F}$ y $\cal{G}$ como categorías sobre~$\cal{E}$); pero, por
supuesto, como en cualquier otra categoría~$\cal{C}_{/S}$, se cometerá
habitualmente el abuso de lenguaje consistente en identificar los
$\cal{E}$-funtores (en el sentido precisado más arriba) con flechas de una
categoría~$\Cat_{/\cal{E}}$.

Denotaremos por
$$
\Hom_{\cal{E}}(\cal{F},\cal{G})
$$
el conjunto de los $\cal{E}$-funtores de~$\cal{F}$ en~$\cal{G}$. Por supuesto,
un compuesto de $\cal{E}$-funtores es un $\cal{E}$-funtor (la composición en
cuestión corresponde, por definición, a la composición de las flechas
de~$\Cat_{/\cal{E}}$).

Consideremos ahora dos $\cal{E}$-funtores
$$
f,g\colon \cal{F}\to\cal{G}
$$
y una transformación natural de funtores:
$$
u\colon f\to g
$$
Se dice que $u$ es un $\cal{E}$-\emph{homomorfismo}, o un
«\emph{homomorfismo de $\cal{E}$-funtores}», si para todo
$\xi\in\Ob(\cal{F})$ se tiene
$$
q(u(\xi))=\id_{p(\xi)} \quoi\text{,}
$$
en palabras: poniendo
\ifthenelse{\boolean{orig}}
{$S=p(\xi)=qf(\xi)=qg(\xi)\in\Ob\cal{E}$,}
{$S=p(\xi)=qf(\xi)=qg(\xi)\in\Ob(\cal{E})$,}
el morfismo
$$
u(\xi)\colon f(\xi)\to g(\xi)
$$
de~$\cal{G}$ es un $\id_S$-morfismo. (En general, para todo morfismo
$\alpha\colon T \to S$ de~$\cal{E}$ y toda categoría~$\cal{G}$ sobre~$\cal{E}$,
un morfismo~$v$ de~$\cal{G}$ se llama $\alpha$-\emph{morfismo} si
$q(v)=\alpha$, donde $q$ denota el funtor de proyección
$\cal{G}\to \cal{E}$). Si se tiene un tercer $\cal{E}$-funtor
$h\colon \cal{F}\to \cal{G}$ y un $\cal{E}$-homomorfismo $v\colon g\to h$,
entonces $vu$ es también un $\cal{E}$-homomorfismo. Así,
\marginpar{150}
\emph{los $\cal{E}$-funtores de~$\cal{F}$ en~$\cal{G}$ y los
$\cal{E}$-homomorfismos entre ellos forman una subcategoría de la categoría
$\SheafHom(\cal{F},\cal{G})$ de todos los funtores de~$\cal{F}$ en~$\cal{G}$,
que se llamará categoría de los $\cal{E}$-funtores de~$\cal{F}$ en~$\cal{G}$ y
se denotará por}
$$
\SheafHom_{\cal{E}/-}(\cal{F},\cal{G})
$$
\label{indnot:fi}\oldindexnot{$\SheafHom_{\cal{E}/-}(\cal{F},\cal{G})$|hyperpage}%
Es también la subcategoría núcleo del par de funtores
\begingroup
\shorthandoff{<>}% Spanish babel otherwise treats Xy-pic's bend delimiters as shorthands.
$$
\xymatrix@C=.5cm{R,S\colon\SheafHom(\cal{F},\cal{G})
\ar@<2pt>[r]\ar@<-2pt>[r]&\SheafHom(\cal{F},\cal{E})}\quoi\text{,}
$$
\endgroup
donde $R$ es el funtor constante definido por el objeto~$p$ de
$\SheafHom(\cal{F},\cal{E})$, y $S$ es el funtor $f\mto q\circ f$ definido por
$q\colon\cal{G}\to\cal{E}$.

Para terminar estas generalidades, queda definir las operaciones naturales de
composición entre las categorías $\SheafHom_{\cal{E}/-}(\cal{F},\cal{G})$
mediante la composición de $\cal{E}$-funtores. En otras palabras, se quiere
definir un
«funtor de composición»:
\begin{equation*}
\label{eq:VI.2.i}
\tag{i} {\SheafHom_{\cal{E}/-}(\cal{F},\cal{G})\times
\SheafHom_{\cal{E}/-}(\cal{G},\cal{H})\to
\SheafHom_{\cal{E}/-}(\cal{F},\cal{H})}
\end{equation*}
cuando $\cal{F}$, $\cal{G}$ y $\cal{H}$ son tres categorías sobre~$\cal{E}$,
de modo que este funtor induzca, en los objetos, la aplicación de composición
$(f,g)\mto gf$ de los $\cal{E}$-funtores
$f\colon\cal{F}\to\cal{G}$ y $g\colon \cal{G}\to\cal{H}$. Para ello,
recordemos que se define un funtor canónico:
\begin{equation*}
\label{eq:VI.2.ii}
\tag{ii}
{\SheafHom(\cal{F},\cal{G})\times\SheafHom(\cal{G},\cal{H})\to
\SheafHom(\cal{F},\cal{H})}
\end{equation*}
que, en los objetos, no es sino la aplicación de composición $(f,g)\mto gf$ de
funtores y que transforma una flecha~$(u,v)$, donde
$$
u\colon f\to f'\quoi,\quad v\colon g\to g'
$$
son flechas de~$\SheafHom(\cal{F},\cal{G})$ y de
$\SheafHom(\cal{G},\cal{H})$, respectivamente, en la flecha
$$
v*u\colon gf\to g'f'
$$
\label{indnot:fj}\oldindexnot{$v*u$|hyperpage}%
definida por la relación:
\marginpar{151}
$$
v*u(\xi)=v(f'(\xi))\cdot g(u(\xi))=g'(u(\xi))\cdot v(f(\xi))
$$
Es bien sabido que así se obtiene efectivamente una transformación natural
de~$gf$ en~$g'f'$ y que (para $f,g$ y $u,v$ variables) se obtiene de este modo
un funtor~\eqref{eq:VI.2.ii}, es decir, que se tiene
\begin{equation*}
\label{eq:VI.2.I}
\tag{I} {\id_g*\id_f=\id_{gf}\quoi,}
\end{equation*}
\begin{equation*}
\label{eq:VI.2.II}
\tag{II} {(v'*u')\circ(v*u)=(v'\circ v)*(u'\circ u)}
\end{equation*}
Recordemos también que existe una fórmula de asociatividad para las
composiciones canónicas~\eqref{eq:VI.2.ii}, expresada, por una parte, por la
asociatividad $(hg)f=h(gf)$ de la composición de funtores y, por otra, por la
fórmula
\begin{equation*}
\label{eq:VI.2.III}
\tag{III} {(w*v)*u=w*(v*u)}
\end{equation*}
para las composiciones horizontales de transformaciones naturales (donde
$u\colon f\to f'$ y $v\colon g\to g'$ son como antes, y se supone dada además
una transformación natural $w\colon h\to h'$ de los funtores
$h,h'\colon \cal{H}\to\cal{K}$). Afirmamos ahora que
\emph{cuando~$\cal{F}$ y $\cal{G}$ son $\cal{E}$-categorías, el funtor de
composición canónico}~\eqref{eq:VI.2.ii}
\emph{induce un funtor}~\eqref{eq:VI.2.i}. Como ya se sabe que el compuesto de
dos $\cal{E}$-funtores es un $\cal{E}$-funtor, esto equivale a decir que
\emph{cuando $u\colon f\to f'$ y $v\colon g\to g'$ son homomorfismos de
$\cal{E}$-funtores, entonces $v*u\colon gf\to g'f'$ es también un homomorfismo
de $\cal{E}$-funtores.} Esto resulta, en efecto, inmediatamente de las
definiciones. Como las operaciones de composición~\eqref{eq:VI.2.i} son
inducidas por las operaciones de composición~\eqref{eq:VI.2.ii}, satisfacen la
misma propiedad de
asociatividad, expresada también por las fórmulas $(hg)f=h(gf)$ y
$(w*v)*u=w*(v*u)$ para $\cal{E}$-funtores y $\cal{E}$-homomorfismos de
$\cal{E}$-funtores.

Para completar el formulario~\eqref{eq:VI.2.I}, \eqref{eq:VI.2.II},
\eqref{eq:VI.2.III}, recordemos también las fórmulas:
\begin{equation*}
\label{eq:VI.2.IV}
\tag{IV} {v*\id_{\cal{F}}=v\quad\text{y}\quad\id_{\cal{G}}*u=u}\quoi,
\end{equation*}
donde
\marginpar{152}
para abreviar se escribe $v*f$ o
\ifthenelse{\boolean{orig}}{$u*g$}{$g*u$} en lugar de~$v*u$, cuando $u$,
respectivamente $v$, es el automorfismo identidad de~$f$, respectivamente
de~$g$.

De la definición de las operaciones de composición~\eqref{eq:VI.2.i} se sigue
que
$\SheafHom_{\cal{E}/-}(\cal{F},\cal{G})$ \emph{es un funtor en
\ifthenelse{\boolean{orig}}
{$\cal{E}$,~$\cal{G}$,}
{$\cal{F}$,~$\cal{G}$,}
de la categoría producto ${\Cat_{/\cal{E}}}^\circ\times\Cat_{/\cal{E}}$ en la
categoría~$\Cat$}. En efecto, si $g\colon \cal{G}\to\cal{G}_1$ es un
$\cal{E}$-funtor, es decir, un objeto de
$\SheafHom_{\cal{E}/-}(\cal{G},\cal{G}_1)$, entonces, haciendo
$\cal{H}=\cal{G}_1$ en~\eqref{eq:VI.2.i}, le corresponde un funtor
$$
g_*\colon\SheafHom_{\cal{E}/-}(\cal{F},\cal{G})\to
\SheafHom_{\cal{E}/-}(\cal{F},\cal{G}_1)
$$
De manera análoga, para un $\cal{E}$-funtor
$f\colon \cal{F}_1\to \cal{F}$ se define un funtor
$$
f^*:\SheafHom_{\cal{E}/-}(\cal{F},\cal{G})\to
\SheafHom_{\cal{E}/-}(\cal{F}_1,\cal{G})
$$
Para abreviar, estos funtores se denotan también por los símbolos
$f\mto g\circ f$ y $g\mto g\circ f$, respectivamente (que, en realidad,
designan únicamente las aplicaciones correspondientes en los conjuntos de
objetos). De la propiedad de asociatividad indicada más arriba se sigue que,
de este modo, se obtiene efectivamente el funtor anunciado
$\Cat_{/\cal{E}}^\circ\times\Cat_{/\cal{E}}\to\Cat$.
